%% S6: every finite prefix of t^inf is computed by a finite S-model; realizable prefixes are decidable.
%% DRAFT 2026-11-23 (skinock).  Neutral amsart shell, as S1--S5.
%% Sources: ~/Code/ski/notes/S4-bounded-exposure.md, sections R1--R8 and the three Wizard reviews.
%% Paper-to-Lean map (lean-ski; all in #audit_clean, propext + Quot.sound only):
%%   lem:kinds ........ PsiSkel.one_iff, one_item_of_fold, fold_two_items        (branch psi-skeleton)
%%   lem:productive ... BExp.skeleton_r1, BExp.r1_le, HnfSound.wt_le_pot           (PR #89)
%%   lem:burst ........ BExp.iterS_run, BExp.fib_dead                              (PR #89)
%%   lem:stable ....... BExp.stable_unique, stable_steps                           (PR #89)
%%   lem:psiskel ...... GCert.psi_skeleton (branch g-certificate, d1c0143): ranks checked by decide +kernel,
%%                      soundness round_inv, headChain_penult, phase; PsiSkel.root_head_not_comb, last_mem_Lz
%%   thm:index ........ pencil. Phi.theorem_phi is an analogue for the hnf observable (grade 10,
%%                      only towers normalized); S6's relation does not imply its hypothesis.
%%   lem:decide, proof of thm:main ... pencil
%% TODO before release: version/DOI of lean-ski (needs psi-skeleton + g-certificate merged); Grammarian, Wizard.
\documentclass[11pt]{amsart}
\usepackage[T1]{fontenc}
\usepackage{lmodern}
\usepackage{amsmath,amssymb,amsthm}
\usepackage[numbers,sort]{natbib}
\usepackage[hidelinks]{hyperref}

\newtheorem{theorem}{Theorem}
\newtheorem{lemma}[theorem]{Lemma}
\newtheorem{corollary}[theorem]{Corollary}
\newtheorem{conjecture}[theorem]{Conjecture}
\theoremstyle{definition}
\newtheorem{definition}[theorem]{Definition}

\newcommand{\Ter}{\mathcal S}
\newcommand{\hole}{\bullet}
\newcommand{\B}{\mathbf B}
\newcommand{\hred}{\twoheadrightarrow_h}
\newcommand{\conv}{=_\beta}

\title{Prefix Regularity for the Combinator $S$}
\author{Claude Opus \& Fable with N.\ E.\ Davis}
\address{ZeroAttest, Champaign, Illinois, USA}
\email{NDavis@ZeroAttest.com}
\urladdr{https://orcid.org/0000-0002-1131-8319}
\date{\today}

\begin{document}
\begin{abstract}
Every closed term $t$ built from the combinator $S$ has an infinitary normal form $t^\infty$. We show
that for every depth $d$ a finite model of the $S$-law computes the depth-$d$ prefix of $t^\infty$: the
fold of $t$ into the model determines it. The models are computable from $d$. Consequently, for every
finite tree $P$ the terms whose prefix is $P$ form a regular tree language, and it is decidable whether
some closed $S$-term has an infinitary normal form beginning with $P$. The proof rests on two
bounds on head reduction, one of which uses a finite computation on a 38-element table.
\end{abstract}

\keywords{combinatory logic, combinator $S$, infinitary normal forms, tree automata, finite models}
\maketitle

\section{Introduction}

$\Ter$ is the set of closed terms built from $S$ by application, with $S\,x\,y\,z\to x\,z\,(y\,z)$ at
any position; $\conv$ is convertibility. Each $t\in\Ter$ has a unique infinitary normal form
$t^\infty$~\citep{waldmann-thesis,s3}, and its prefix of each finite depth is computable~\citep{s3}.
A 38-element model of the $S$-law computes the prefix of depth one~\citep{s1}, and a 373-element model
computes the prefix of depth two~\citep{s5}. We show that such models exist at every depth.

\begin{theorem}[Prefix Regularity]\label{thm:main}
For every $d\ge0$ there are a finite $S$-model $M_d$, generated by its element $s$, and a map $o_d$
from $M_d$ to trees such that the depth-$d$ prefix of $t$ is $o_d(\bar t)$ for every $t\in\Ter$, where
$\bar t$ is the fold of $t$ into $M_d$. The table of $M_d$ and the map $o_d$ are computable from $d$.
\end{theorem}

\begin{corollary}\label{cor:realize}
For every $d$ and every finite tree $P$, the set of $t\in\Ter$ whose depth-$d$ prefix is $P$ is a
regular tree language, and it is decidable whether this set is empty. In particular, it is
decidable whether some closed $S$-term has an infinitary normal form whose depth-$d$ prefix is $P$.
\end{corollary}

\begin{proof}
The fold into $M_d$ is a deterministic bottom-up tree automaton with states $M_d$, and the set is
the preimage of $o_d^{-1}(P)$. Since $M_d$ is generated by $s$, every element is the fold of a term, so
the set is nonempty iff $P\in o_d(M_d)$. The depth-$d$ prefix of $t^\infty$ is $\pi_d(t)$ with the head
$S$ of each hnf cut~\citep{s5}, a computable map of $\pi_d(t)$, so the prefixes of infinitary normal
forms that occur are the finitely many cut images of $o_d(M_d)$.
\end{proof}

A prefix may depend on material arbitrarily deep in $t$ (Section~\ref{sec:close}), so no bounded
inspection of $t$ suffices; the finite models absorb the unbounded part by an eventual periodicity of
one-item chains $S\,(S\,(\cdots(S\,u)))$.

\section{Frames, prefixes and observables}

Every $t\in\Ter$ is $S\,u_1\cdots u_m$ for a unique $m\ge0$; this is its \emph{frame}, the $u_i$ are its
\emph{items}, and $t$ is a \emph{head normal form} (hnf) iff $m\le2$. If $m\ge3$, the \emph{head step}
is $t\to u_1\,u_3\,(u_2\,u_3)\,u_4\cdots u_m$; it is \emph{free} if $u_1$ is $S$ or of the form
$S\,v$, and \emph{non-free} otherwise. Head steps define the \emph{head run} of $t$, and
$\hred$ is its reflexive--transitive closure. A frame with $m\ge3$ is \emph{stable} if
$S\,u_1\cdots u_{m-1}$ has no hnf; we write it $f\,w$ with $f=S\,u_1\cdots u_{m-1}$ and $w=u_m$.

An \emph{$S$-model} is a set $M$ with a binary operation and an element $s$ such that
$s\,a\,b\,c=a\,c\,(b\,c)$; its \emph{fold} sends $S$ to $s$ and $t\,u$ to $\bar t\,\bar u$, and
$\bar t=\bar u$ whenever $t\conv u$. Let $\B$ be the 38-element model of~\citep{s1}, whose elements
are the \emph{blocks}, with $\beta$ its operation and $0$ the block of $S$. It has an element $\bot$
with $\bot\,a=\bot$, and a term has an hnf iff its block is not $\bot$~\citep{s1}. The head run of every
term reaches an hnf or a stable frame~\citep{s1,s3}. Once a run is at a stable frame $f\,w$, every
later head step is a head step of $f$: $w$ is kept, and $f$ is replaced by a reduct, so $\bar f$ is unchanged in every model. Throughout,
$M$ is an $S$-model with a homomorphism to $\B$ that fixes $s$, so $\bar t$ determines the block of $t$.

\begin{definition}\label{def:obs}
The \emph{observable} of $t$ over $M$ is $O(t)=(m,\bar u_1,\dots,\bar u_m)$ if the head run of $t$
reaches the hnf $S\,u_1\cdots u_m$, and $O(t)=(\bot,\bar f,\bar w)$ if it reaches the stable frame
$f\,w$ first. For $t$ without an hnf we write $\Psi(t)=(\bar f,\bar w)$.
\end{definition}

The prefixes are those of~\citep[Def.~2]{s5}, written $\pi_d$, with the head $S$ of an hnf shown.
Trees are built from $S$, a hole $\hole$ and binary nodes; $\pi_0(t)=\pi_{-1}(t)=\hole$, and
$\pi_{d+1}(t)$ is $S$, $S\,\pi_d(a)$, $S\,\pi_{d-1}(a)\,\pi_d(b)$ or $\pi_d(f)\,\pi_d(w)$, according as
the head run of $t$ reaches $S$, $S\,a$, $S\,a\,b$ or the stable frame $f\,w$ first.

\begin{lemma}\label{lem:invariant}
If $t\conv t'$, then $O(t)=O(t')$.
\end{lemma}

\begin{proof}
Let $T$ be the hnf or first stable frame on the head run of $t$. If $T$ is an hnf, its reducts are frames of the same length whose items are reducts of those of $T$;
if $T=f\,w$ is stable, its reducts are $f^*\,w^*$ with $f\twoheadrightarrow f^*$ and
$w\twoheadrightarrow w^*$, again stable, because $f^*$ keeps the block $\bot$. By Church--Rosser, $t'$ and $T$ have a common reduct $T^*$. By
standardization~\citep{terese}, $t'\hred T'\twoheadrightarrow T^*$ with only internal steps
after $T'$; these keep the frame length and the block of every prefix, so $T'$ is an hnf or a stable
frame of the same length as $T^*$, and its items reduce to those of $T^*$. Hence corresponding items of $T$ and $T'$ (for stable frames, $f,f'$ and $w,w'$, where
$T'=f'\,w'$) are convertible and have equal values in $M$. If $T'$ is stable, the first stable frame
of $t'$ precedes it on the run and, by the remark above, has the same value pair.
\end{proof}

\section{Two bounds on head runs}

\begin{lemma}[kinds]\label{lem:kinds}
The block of $t$ determines whether $t$ is $S$, of the form $S\,v$, or has at least two items:
for all blocks $a,b$, $\beta(a,b)$ lies in the set $\mathrm{ONE}$ of the blocks of the terms $S\,v$ iff $a=0$, and $0$ is the
block of $S$ alone.
\end{lemma}

\begin{proof}
A finite computation on $\beta$, with~\citep{s1} for block $0$.
\end{proof}

\begin{lemma}[productive runs]\label{lem:productive}
There is a map $\varpi$ from the blocks other than $\bot$ to $\{0,\dots,10\}$ such that the head run of
a term $t$ with an hnf has at most $\varpi(\bar t)$ non-free steps.
\end{lemma}

\begin{proof}
Give a frame the weight obtained by reading its items' blocks through a table of potentials on $\B$.
A finite computation on $\beta$ shows that the weight of a term with an hnf is at most $\varpi$ of its
block, that no head step raises it, and that every step whose first item has a block other than $0$
and $17$ lowers it. Blocks $0$ and $17$ are those of $S$ and of the terms $S^k(S\,S)$, $k\ge4$, so a
non-free step lowers the weight by Lemma~\ref{lem:kinds}.
\end{proof}

For terms $x,y$ let $(P_0,Q_0)=(x,y)$ and $(P_{j+1},Q_{j+1})=(Q_j,P_j\,Q_j)$.

\begin{lemma}[bursts]\label{lem:burst}
Let $k\ge0$. Then $S^k u\,x\,y\hred u\,P_k\,Q_k$ in $k$ head steps, the $j$-th with first item
$S^{k-j}u$. For $k\ge8$, $P_k$ and $Q_k$ have no hnf; for $k\ge10$, $P_k=P_{k-2}\,P_{k-1}$ and
$Q_k=P_{k-1}\,P_k$ are stable frames.
\end{lemma}

\begin{proof}
$S\,(S^{k-1}u)\,x\,y\to S^{k-1}u\,y\,(x\,y)$; induct. On blocks, $(a,b)\mapsto(b,\beta(a,b))$ maps
every pair to $(\bot,\bot)$ within eight steps, a finite computation. The recurrence gives
$P_k=Q_{k-1}=P_{k-2}\,Q_{k-2}=P_{k-2}\,P_{k-1}$, and $f\,w$ with $f$ lacking an hnf is a stable frame.
\end{proof}

A \emph{root frame} is a frame $S\,x\,y\,z$ of three items with no hnf.

\begin{lemma}[$\Psi$-skeleton]\label{lem:psiskel}
The head run of a term without an hnf has at most $77$ non-free steps before its first stable frame.
\end{lemma}

\begin{proof}
If the frame of $t$ has $m\ge4$ items and is not stable, its prefix has an hnf, and the run of $t$
follows that of the prefix (at most $10$ non-free steps, Lemma~\ref{lem:productive}) to a root frame.
From a root frame $S\,x\,y\,z$ one step gives $h\,w$ with $h=x\,z$, $w=y\,z$. If $h$ has no hnf, $h\,w$
is stable. Otherwise $h$ runs to an hnf with exactly two items, $S\,p\,q$, since fewer would make
$h\,w$ an hnf; the frames in between have prefixes with an hnf, so none is stable, and $S\,p\,q\,w$ is
the next root frame. Call this a \emph{round}; it has at most $1+\varpi(\bar h)$ non-free steps. We
collect what the blocks force.
\begin{enumerate}
\item The block of $x$ is not that of $S$ or of a term $S^j(S\,S)$, since $S\,x\,y\,z$ would then have
  an hnf (Lemma~\ref{lem:burst}).
\item If $x=S\,a$, then $h=S\,a\,z$, so $(p,q)=(a,z)$ and the round has no non-free step.
\item Otherwise $h$ has at least three items, and its run ends with a free step
  $S\,S\,y'\,z'\to S\,z'\,(y'\,z')$, because a head step shortens a frame by at most one item and only
  a step from three items whose first item is $S$ reaches two. So $p=z'$ and $q=y'\,p$. The last item of a frame
  changes only at a step from three items, $[x',a,b]\mapsto[\dots,b,a\,b]$, so $\bar p$ lies in the
  closure $L(\bar z)$ of $\{\bar z\}$ under $e\mapsto\beta(c,e)$, $c\in\B$.
\end{enumerate}
\sloppy
Let $G_1$ be the graph on the triples $(\bar x,\bar y,\bar z)$ allowed by (1) with
$\beta(\beta(\beta(0,\bar x),\bar y),\bar z)=\bot$, with an edge to every triple that (2) or (3)
permits for the next root frame. Successive root frames of a run form a path in $G_1$. On a set $F$ of $190$ triples, all with $\bar z\in\{31,\bot\}$, a refinement $G_2$ records $x$ as
$c\,e$ and follows the round through the blocks of the frame prefix of $h$, of its last two items
and, once the last item has changed, of its argument part, within a set of round states closed
under the moves that $\beta$ forces; like $G_1$, it over-approximates the run. Weight a node by
$1+\varpi(\bar h)$ if $x$ has at least two items, reading $\varpi(\bot)$ as $0$, and by $0$
otherwise; this bounds the non-free steps of its round, the last included. A finite computation
gives ranks, at most $67$, on the nodes of $G_2$ over $F$ and of $G_1$ elsewhere, each at least the
weight of its node, that drop by at least the weight of the source along every edge of $G_2$ and
every edge of $G_1$ leaving a triple outside $F$. So the weights along a run sum to at most $67$, and the run has at most
$10+67$ non-free steps before its first stable frame.
\end{proof}

\section{Finite index}

Fix $M$ finite. Let $\sigma(a)=s\,a$ and $\phi(a,b)=(b,a\,b)$ on $M$ and $M^2$. Since $M$ is finite,
there are $N\ge10$ and $p\ge1$ with $\sigma^j=\sigma^{j'}$ and $\phi^j=\phi^{j'}$ whenever
$j,j'\ge N-2$ and $j\equiv j'\pmod p$. Then $\overline{S^k u}=\sigma^k(\bar u)$, and by
Lemma~\ref{lem:burst}, for $k\ge10$, $\Psi(P_k)$ and $\Psi(Q_k)$ are read off
$\phi^{k-2}(\bar x,\bar y)$, $\phi^{k-1}(\bar x,\bar y)$ and $\phi^k(\bar x,\bar y)$.

By Lemma~\ref{lem:kinds}, every term is uniquely $S^k v$ with $v=S$ or $v$ of at least two items; this is
its \emph{core form}. Write $k\sim k'$ if $k=k'$, or $k,k'\ge N$ and $k\equiv k'\pmod p$. Define
$v\approx_{-1}v'$ iff $\bar v=\bar v'$ and, if $v$ has no hnf, $\Psi(v)=\Psi(v')$. For $g\ge0$,
$v\approx_g v'$ iff $v\approx_{-1}v'$ and, if $v$ has an hnf $S\,a\,b$ (then so does $v'$, say
$S\,a'\,b'$, by equal blocks and Lemma~\ref{lem:kinds}), $a\simeq_{g-1}a'$ and $b\simeq_{g-1}b'$.
Finally $u\simeq_g u'$ iff their core forms $(k,v)$, $(k',v')$ satisfy $k\sim k'$ and $v\approx_g v'$.
By induction on $g$, $\simeq_g$ is an equivalence of finite index, and $u\simeq_g u'$ implies
$\bar u=\bar u'$.

\begin{theorem}\label{thm:index}
If $t\simeq_{77}t'$, then $O(C[t])=O(C[t'])$ for every context $C$.
\end{theorem}

\begin{proof}
Let $\rho(P)$ be the number of non-free steps on the head run of $P$ before its hnf or first stable
frame; $\rho\le77$ by Lemmas~\ref{lem:productive} and~\ref{lem:psiskel}. A \emph{configuration} is
$(P,P',\Pi)$ where $P'$ arises from $P$ by replacing disjoint occurrences of $u$ by $u'$ for the pairs
$(u,u')\in\Pi$, each of one of two kinds: (A) $u,u'$ without hnf and $u\approx_{-1}u'$; or (C) $u\simeq_g
u'$ with a recorded grade $g\ge-1$, where the core forms satisfy $k\ne k'$, or $v$ has at least two
items and an hnf. Every other $\simeq$-related pair is an identity or $S^k$ applied to an (A) pair, and
is recorded as such. All pairs have equal values, so $P$ and $P'$ have the same block, and so do
corresponding prefixes. A pair is \emph{at the head} if it is a prefix of the spine of $P$. Invariant:
(I1) a (C) pair at the head with $k=k'=0$ has $g\ge\rho(P)$; (I2) every other (C) pair has
$g\ge\rho(P)-1$. The start $(C[t],C[t'],\{(t,t')\})$, with grade $77$, satisfies it. Throughout, $P$ and $P'$ advance along their own head runs without passing an hnf or a stable frame, so $O$ is kept;
we induct on the number of head steps left to $P$, and then on whether a pair is at the head. The cases are tested in the order given, so in
(H1)--(H3) and (N) $P$ is neither an hnf nor stable.

\emph{(T) $P$ is an hnf or a stable frame, or an (A) pair is all of $P$.} In the last case
$O(P)=(\bot,\Psi(P))=(\bot,\Psi(P'))=O(P')$. If no pair is at the head, every pair lies inside an
item, so $P'$ has the same frame with items of equal values, and is an hnf, respectively a stable
frame. If $P$ is stable, a pair at the head lies in $f$ or is all of $P$; in the first case
$P'=f'\,w'$ with $\bar f'=\bar f=\bot$ in $\B$ and $\bar w'=\bar w$, in the second it is an (A) pair and
$\Psi$ agrees. If $P$ is an hnf, a pair at the head is $S^k v$ with at most one argument, whose
items are $S^{k-1}v$ on one side and $S^{k'-1}v'$ on the other with equal values; or it is all of $P$, a pair $(v,v')$ with $k=k'=0$,
where $g\ge\rho=0$ by (I1) and $v\approx_0v'$ gives items of equal values.

\emph{(H1) A pair at the head with $k=k'\ge1$.} $S^kv=S\,(S^{k-1}v)$: record $(S^{k-1}v,S^{k-1}v')$
instead, now an item. No step is taken and no pair is left at the head.

\emph{(H2) A pair at the head with $k\ne k'$ and at least two arguments.} By Lemma~\ref{lem:burst},
$S^kv\,x\,y\,\vec r\hred v\,P_k\,Q_k\,\vec r$ through frames whose prefix keeps its block, so no hnf or
stable frame is passed; likewise on the other side. As $\bar x=\bar x'$, $\bar y=\bar y'$ and
$k,k'\ge N$, the pairs $(P_k,P'_{k'})$ and $(Q_k,Q'_{k'})$ are of kind (A). The last step has first
item $v$; if it is non-free, $\rho$ drops by one and the new head pair $(v,v')$ satisfies (I1) by
(I2); if $v=S$, then $v'=S$.

\emph{(H3) A (C) pair $(v,v')$ at the head with $k=k'=0$.} Run $v$ to its hnf $S\,a\,b$ inside $P$, and
$v'$ to $S\,a'\,b'$; again no hnf or stable frame of $P$ is passed. If no argument follows, this is
(T). Otherwise replace the pair by $(a,a')$ and $(b,b')$ with grade $g-1\ge\rho-1$ by (I1); they are
not at the head. A pair of grade $-1$ never reaches this case: by (I2), $\rho=0$, while a core of at
least two items reaches the head only by a non-free step.

\emph{(N) No pair at the head, $P$ not an hnf and not stable.} Take the head step on both sides; it
copies $u_3$, keeping occurrences disjoint. A core $v$ of at least two items reaches the head only if
the first item was $v\,\vec c$, which has at least two items, so the step is non-free, $\rho$ drops,
and (I2) gives (I1). Other pairs keep (I2), since $\rho$ never grows.
\end{proof}

\section{The tower of models}

\begin{lemma}\label{lem:decide}
If $M$ is finite with a given table, the syntactic congruence $\equiv_O$ of $O$ over $M$, given by $t\equiv_O t'$ iff
$O(C[t])=O(C[t'])$ for all contexts $C$, is decidable, and $\Ter/{\equiv_O}$ is a finite $S$-model
generated by $[S]$ whose table is computable.
\end{lemma}

\begin{proof}
$O$ is computable: the block decides whether the run stops at an hnf or at a stable frame, and the run
reaches one of them in finitely many steps. By Theorem~\ref{thm:index}, $\equiv_O$ has finite index at most $n$, where $n$ bounds the
index of $\simeq_{77}$ and is computable from $|M|$, $N$ and $p$. Let $R$ be the terms of depth at
most $n$; they meet every class, because the classes met by terms of depth at most $i$ grow
strictly with $i$ until they are closed under application. Let $E_0$ be the kernel of $O$ and $E_{i+1}$ the pairs of $E_i$ with
$(t\,r,t'\,r),(r\,t,r\,t')\in E_i$ for all $r\in R$. Each $E_i$ contains $\equiv_O$; when $E_{i+1}=E_i$,
$E_i$ is a congruence inside the kernel of $O$, as every term is $\equiv_O$ to one in $R$, so
$E_i={\equiv_O}$. The index of $E_i$ grows at each strict step and is at most $n$, so $E_n={\equiv_O}$, and $E_n$
is decided by finitely many evaluations of $O$. By Lemma~\ref{lem:invariant}, $\equiv_O$ contains
$\conv$, so the quotient satisfies the $S$-law. Its table is found by closing $\{S\}$ under
application modulo $\equiv_O$.
\end{proof}

\begin{proof}[Proof of Theorem~\ref{thm:main}]
We build $M_d$, $o_d$ and homomorphisms $h_d\colon M_d\to M_{d-1}$ fixing $s$. Let $M_0$ be the
one-element model with $o_0=\hole$, and $M_1=\B$, whose block determines $\pi_1$~\citep{s1}. Given
$M_{d-1}$ and $M_d$ for some $d\ge1$, let $O$ be the observable over $M_d$. By the definition of
$\pi_{d+1}$, $\pi_{d+1}(t)$ is $S$, $S\,o_d(a)$, $S\,o_{d-1}(h_d\,a)\,o_d(b)$ or $o_d(f)\,o_d(w)$ when
$O(t)$ is $(0)$, $(1,a)$, $(2,a,b)$ or $(\bot,f,w)$, since $h_d$ composed with the fold into $M_d$ is
the fold into $M_{d-1}$. Let $M_{d+1}=\Ter/{\equiv_O}$ (Lemma~\ref{lem:decide}), $o_{d+1}([t])$ the tree $\pi_{d+1}(t)$, and $h_{d+1}([t])=\bar t\in M_d$, which is well defined because $O(t)$ determines $\bar t$: it is
$s\,\bar u_1\cdots\bar u_m$ or $\bar f\,\bar w$. Then $h_{d+1}$ is a homomorphism, and $M_{d+1}$ maps to
$\B$ through $M_d$.
\end{proof}

\section{Closing}\label{sec:close}

No fixed inspection depth suffices. Let $c_1=S\,(S\,S)$ and $c_2=S\,S\,S$. For $k\ge8$ and all $x,y$,
the terms $t_i=S^k(S\,S\,c_i)$ satisfy $t_i\,x\,y\hred P\,Q\,(c_i\,P\,Q)$ with $(P,Q)=(P_k,Q_k)$, both without an hnf,
and $\overline{c_1\,P\,Q}=31\ne\bot=\overline{c_2\,P\,Q}$ in $\B$, while $S^j(S\,S\,c_1)$ and
$S^j(S\,S\,c_2)$ have equal blocks for every $j$. So the depth-$2$ prefix of $t_i\,x\,y$ depends on $c_i$,
which lies $k+1$ levels down in $t_i$, and the block-labelled hnf unfoldings of $t_1$ and $t_2$ agree to depth $k$. Theorem~\ref{thm:index} absorbs this through the periodicity of $k$.

The models grow quickly, and the bound of
Theorem~\ref{thm:index} is far from sharp. Convertible terms have equal infinitary normal forms, and
Theorem~\ref{thm:main} separates terms with different ones by finite models. By McKinsey and
Evans~\citep{mckinsey,evans}, convertibility would therefore be decidable if the converse held.

\begin{conjecture}
If $t^\infty=u^\infty$, then $t\conv u$.
\end{conjecture}

\section*{Data availability}
Lemmas~\ref{lem:kinds}, \ref{lem:productive}, \ref{lem:burst} and~\ref{lem:psiskel}, the latter including the soundness of the round abstractions behind $G_1$ and $G_2$ and the check
of the ranks, are proved in the
Lean~4~\citep{lean4} development \texttt{lean-ski}~\citep{leanski}, using only the axioms
\texttt{propext} and \texttt{Quot.sound}. The same development proves an analogue of Theorem~\ref{thm:index} for the
head-normal-form part of $O$, with grade $10$ and only the towers $S^k(S\,S)$ normalized. Facts about the strongly connected components of $G_1$ and $G_2$ are computed only and are not
used. The remaining arguments are proved here.

\bibliographystyle{plainnat}
\bibliography{refs}
\end{document}
